\section{Introduction}

What happens when a university rewrites its intellectual-property policy? Research universities generally maintain formal intellectual-property (IP) policies that specify who owns an invention, what faculty must disclose and when, how licensing income is shared, how conflicts of interest are handled, and what inventors can expect from the technology-transfer office (TTO). Universities do not leave these documents unchanged forever. From time to time, often after many years, they rewrite them in response to new law, new leadership, changing governance priorities, or experience with commercialization.

A large literature has studied the institutional arrangements behind university technology transfer. One line of work examines who owns university inventions and how ownership affects incentives to commercialize \citep{HvideJones2018,KenneyPatton2009}. Another asks whether direct financial incentives, especially inventor royalty shares, affect invention and commercialization. The evidence is mixed. \citet{LachSchankerman2008} reported that universities offering inventors larger royalty shares generated more licensing income, and \citet{ArqueCastellsEtAl2016} linked royalty sharing to inventive effort in Portugal and Spain. Using corrected and expanded U.S. policy data, however, \citet{OuelletteTutt2020} showed that the earlier licensing-income relationship was driven by coding errors, and they found no compelling evidence that larger inventor royalty shares significantly affected university patenting or commercialization outcomes. Other studies focus on the organization of technology transfer itself, including TTO staffing, routines, incentives, and relationships with industry \citep{SiegelWaldmanLink2003,MarkmanEtAl2005,DebackereVeugelers2005}.

The literature also points to parts of commercialization that are harder to observe. Faculty decisions to bring inventions into the formal technology-transfer process depend partly on how researchers view the TTO and on the commercialization behavior and norms of those around them \citep{OwenSmithPowell2001,BercovitzFeldman2008}. The TTO sits between faculty, university administration, and potential licensees, and it has objectives and decision processes of its own \citep{JensenThursbyThursby2003}. Because many university inventions are licensed at an early stage, commercialization can also require continued cooperation among inventors, TTO staff, and licensees \citep{JensenThursby2001}. Technology transfer thus depends not only on ownership, financial incentives, and office resources, but also on the institutional setting in which inventors and universities interact.

Part of that governance relationship is written into the IP policy, yet this dimension has been difficult to compare systematically across universities and over time. Ownership rules and royalty shares can be coded fairly directly. How a university presents its relationship with inventors is harder to pin down. Two policies may contain broadly similar rules while presenting them differently: one may emphasize obligations, compliance, and institutional authority, while another places more weight on guidance, collaboration, and assistance. The companion paper \citep{SharmaPCSI2026} develops the Policy Communication Stance Index (PCSI) to measure this dimension. PCSI captures how formal IP-policy text presents rights, obligations, institutional authority, and assistance. Higher values indicate a more supportive stance, and lower values a more restrictive one.\footnote{Following the annotation rubric in \citet{SharmaPCSI2026}, supportive language emphasizes assistance, guidance, encouragement, collaboration, shared interests, or facilitation. Restrictive language emphasizes mandatory assignment, categorical obligations, unilateral institutional authority, prohibitions, sanctions, forfeiture, or compliance. These labels describe how the rules are communicated, not whether the underlying legal rule is more or less generous.} PCSI describes the stance of the text itself. It does not measure whether the underlying rules are generous, whether faculty approve of them, how effectively the TTO operates, or how successful the university is at commercialization.

There are several reasons why this written stance could be related to commercialization. Inventors may use the policy to understand their obligations and the help available to them, and formal procedures can create learning and compliance costs \citep{MoynihanHerdHarvey2015}. Policy may also matter indirectly through the TTO, whose routines, incentives, governance structures, and decision processes shape commercialization \citep{JensenThursbyThursby2003,SiegelWaldmanLink2003,DebackereVeugelers2005,MarkmanEtAl2005}. Or a rewrite may simply be the visible part of a broader organizational change, such as new leadership, new workflows, or a new commercialization strategy, since formal rules can be partly decoupled from what organizations actually do \citep{MeyerRowan1977}. Our data cannot distinguish these mechanisms, but they suggest where to look. If a revision mainly changes whether faculty enter the formal technology-transfer process, invention disclosures are a natural early margin. If the relevant changes occur after disclosure, differences may instead appear farther downstream, in patenting or licensing.

To study these revisions, we built a dataset of dated policy-revision events. We begin with the 518 policy documents assembled in the companion paper, covering 150 U.S. research universities from 1944 to 2025, and identify consecutive policy documents at the same university whose PCSI score changes by more than 0.03, roughly half the cross-document standard deviation. Among universities linked to AUTM data, this produces 126 threshold revisions. After mechanical screens for overlapping events and outcome coverage, we read the remaining predecessor and successor documents side by side and keep revisions whose documents are comparable and whose timing can be established with reasonable confidence. The preferred sample contains 34 revisions at 32 universities. Eight are \emph{upward} revisions, in which the successor policy shifts toward a more supportive communicative stance, and 26 are \emph{downward} revisions, in which it shifts toward a more restrictive one. These labels refer only to the direction of change in PCSI. They do not imply that an upward policy is better or more generous, or that a downward policy is legally harsher. Revisions may also change substantive provisions governing disclosure, assignment, revenue sharing, conflicts of interest, equity, or other parts of the university--inventor relationship, so we treat them as revision episodes and not as changes in wording alone. No step in building the sample uses the values of the outcomes.

We link these revisions to technology-transfer outcomes from the AUTM Licensing Activity Survey: invention disclosures, new patent applications, patents issued, an application-to-disclosure filing margin, and licenses and options executed. Our main design is a stacked event study in which each revising university is compared with contemporaneous universities that had no threshold revision in the same event window \citep{CengizEtAl2019,BakerLarckerWang2022}. Because the five outcomes sit at different points in the commercialization process, we can ask where the clearest divergence after a revision appears.

The clearest pattern is in licensing. After upward revisions, licenses and options executed rise relative to contemporaneous non-revisers; after downward revisions, they fall. The average upward-minus-downward difference is about 0.5 log points over the six years after a revision. We find no detectable difference in licensing trends before the revision, and the estimate is similar in a stricter sample and when individual revisions are removed one at a time. The filing margin and new patent applications move in the same direction but less precisely, patents issued show a weaker difference, and invention disclosures show no corresponding increase. The concentration of the pattern downstream of disclosure fits more naturally with mechanisms operating after disclosure, such as evaluation, marketing, negotiation, continued inventor--TTO interaction, or broader organizational reform, than with a simple story in which more supportive revisions mainly lead more faculty to disclose inventions. The aggregate data cannot tell us which of these mechanisms is responsible.

We do not interpret the result as evidence that more supportive wording causes more licensing. Universities choose when and how to revise their policies, and upward revisions start from predecessor policies with much lower PCSI values than downward revisions, which raises real concerns about selection and mean reversion. The licensing evidence also becomes weaker after adjustment for testing multiple outcomes. We treat the estimates as an outcome-specific pattern around policy-revision episodes, not as a causal effect of communicative stance.

The paper brings ordinary university IP-policy revisions into the empirical technology-transfer literature as dated within-institution events. Earlier studies have compared particular policy terms across universities or examined large changes in invention ownership. We instead study the smaller policy revisions that universities periodically undertake themselves and trace outcomes from disclosure through licensing around those episodes. To do so, we built a dataset of documented, dated revisions; PCSI, developed in the companion paper, provides a consistent way to classify their direction across institutions and over time. The evidence shows that the clearest divergence following revision episodes appears at licensing rather than at disclosure, while also showing why these episodes should not be interpreted as exogenous changes in policy language. The paper thus adds evidence on an understudied margin of university technology-transfer governance: the ordinary policy changes universities make within an otherwise stable system of institutional ownership.

Section~\ref{sec:lit} reviews the related literature. Section~\ref{sec:data} describes the policy and AUTM data and the construction of the revision sample. Section~\ref{sec:strategy} presents the empirical design and inference. Section~\ref{sec:results} reports the results and robustness checks. Section~\ref{sec:discussion} discusses interpretation and limitations, and Section~\ref{sec:conclusion} concludes.
